\uuid{3TSh}
\exo7id{3816}
\titre{exo7 3816}
\auteur{quercia}
\organisation{exo7}
\datecreate{2010-03-11}
\isIndication{false}
\isCorrection{false}
\chapitre{Espace euclidien, espace normé}
\sousChapitre{Problèmes matriciels}
\cpge{mpsi-mp2i.calcul-matriciel-systemes.e, pcsi.determinants.c, ptsi.determinants.a, bcpst1.matrices, mp-mpi.reduction-endomorphismes.b, pc.reduction.a, psi.reduction-endomorphismes.a, pt.reduction-endomorphismes.a, bcpst2.valeurs-propres-vecteurs-propres.a}
\module{Algèbre}
\niveau{L2}
\difficulte{}

\contenu{
\question{
Soit~$M\in\mathcal{M}_n(\C)$ et $\lambda\in\mathrm{spec}(M)$. Montrer que $\Re(\lambda)$
est compris entre la plus grande et la plus petite valeur propre de
$\frac12(M+M^*)$.
}
}
